% ============================================================
% f2014_hunan_science_q07_views.tex
% 2014 湖南卷（理科）第 7 题：石材几何体的三视图
% 来源 PDF：普通高考/2014/2014湖南理.pdf
% 原页：第 1 页（A3 横排）第 1 栏底部，印面 885，第 7 题题干下方
% 结构：正视图为直角边长 6、8 的直角三角形；侧视图为矩形；
%       俯视图为 6×12 的矩形，标注柱长 12。比例 0.4 cm/单位。
% ============================================================
\begin{tikzpicture}[>=Stealth, line width=0.6pt, font=\small, scale=0.72,
  x=1cm, y=1cm, line cap=round, line join=round]
  % 正视图：直角边 6（横）、8（纵）
  \draw (0,0) -- (2.4,0) -- (0,3.2) -- cycle;
  \draw[<->] (0,3.95) -- (2.4,3.95);
  \node at (1.2,4.22) {$6$};
  \draw[<->] (-0.95,0) -- (-0.95,3.2);
  \node at (-1.18,1.6) {$8$};
  \node at (1.2,-0.52) {正视图};
  % 侧视图：12×8
  \draw (4.4,0) rectangle (9.2,3.2);
  \node at (6.8,-0.52) {侧视图};
  % 俯视图：6×12
  \draw (10.4,-1.6) rectangle (12.8,3.2);
  \draw[<->] (13.10,-1.6) -- (13.10,3.2);
  \node at (13.52,0.8) {$12$};
  \node at (11.6,-2.12) {俯视图};
\end{tikzpicture}
